\chapter{Methodology}
\label{methodology}

This section describes the chosen methodologies for the project and the rationale for their selection.
The methodologies chosen for this project are the \ac{CDA} and \ac{RAD} model, which will be described in the following sections.

\section{Conceptual-Deductive Analysis}

The \ac{CDA} is a methodology especially popular within the fields of \ac{IS} and \ac{IT} research,
as it has proven itselg to be especially useful as a design basis for the development of complex computer systems\footcite{wildeMethodenspektrumWirtschaftsinformatikUeberblick}.

For this reason it has been chosen as the main co-methodology of many prototype based projects. \footcite{schreinerGestaltungsorientierterKernOder}
Especially in the context of projects where the requirements and possibilities are not fully known at the beginning of the project,
or the previous work in the field is rapidly changing, needing complete evaluation of the current state of the art.

It is based on the following steps:

\begin{enumerate}
    \item Presentation of basic principles
    \item Mapping of the problem domain
    \item Evaluation of previous work and partial solutions
    \item Design of a solution hypothesis
    \item Implementation of the solution
    \item Evaluation of the solution hypothesis compared to the implementation
\end{enumerate}

For this project the \ac{CDA} methodology will be augmented with the \ac{RAD} model a popular prototyping methodology, which will be described in the next section.
The solution hypothesis will be the basis for the requirements of the entire prototyping process.
The final evaluation of the resulting prototype serves as the test of the solution hypothesis,
making this not only a synergistic combination of methodologies, but also a complete scientific method\footcite{wildeMethodenspektrumWirtschaftsinformatikUeberblick}.


\section{Prototyping}

As described in the section above, the combination of Prototyping and \ac{CDA} has been chosen for their ability to handle complex and uncertain problem domains.
While the \ac{CDA} lays the groundwork of the research, the prototyping will be used to evaluate the solution hypothesis and to develop the final product
and grounds the theoretical work in practical application.

Many prototyping methodologies have been developed over the last 40 years \footcite{gomaaImpactRapidPrototyping1983} as they
break down a large and complex problem into smaller achievable steps, which can be tackled individually\footcite{naumannPrototypingNewParadigm1982}.

For this project the \ac{RAD} model has been chosen, as it is signified both by its ability to handle complex and uncertain problem domains,
but also the reduced overhead in comparison to other prototyping methodologies\footcite{martinRapidApplicationDevelopment1991}.
While other methodologies like the spiral model \footcite{boehmSpiralModelSoftware1988} would provide a more thorough risk analysis,
they would also be too time-consuming for a project of this scope, especially since the utilization of the \ac{CDA} already provides a risk mitigation,
by providing a solid theoretical basis for the project.

The \ac{RAD} model is based on the following steps:

\begin{enumerate}
    \item An overall goal is defined which this entire prototype incrementally should approach.
    \item The project is divided into smaller iterations, which are planned and executed individually.
          \begin{itemize}
              \item Specific goal for this iteration are defined
              \item Implementation is documented and deviations from the plan are noted
              \item The results of the implementation and the original plan are evaluated
          \end{itemize}
    \item The results of the previous iteration are used to plan the next iteration
    \item A final evaluation of the original assumptions, the results of the iterations and the final prototype is done
\end{enumerate}

But as this project aims to integrate the \ac{CDA} and the \ac{RAD} model,
the \ac{RAD} model will be augmenting the Prototyping methodology with the \ac{CDA} as described in the previous section.
The solution hypothesis of the \ac{CDA} will then serve as the basis for the requirements of the entire prototyping process,
and the evaluation of the resulting prototype in turn serves as primary test of the solution hypothesis.

\pagebreak
